%% kapa: exam=1 round=34 year=2023 period=2 form=A subject=appraisal_law subject_ko=감정평가관계법규 seq=7 number=7
%% source: Q-Net 감정평가사 기출문제 (cst00309) artlSeq=5211789 file=raw/1008/5211789/2023년도 제34회 감정평가사 1차시험 2교시.pdf
%% conversion: kordoc 4.24.51 -> tools/gen1.py (auto); review=proofread-sonnet; flags=none
\begin{question}{7}
국토의 계획 및 이용에 관한 법령상 도시·군계획시설부지의 매수청구에 관한 설명으로 옳지 않은 것은? (단, 조례는 고려하지 않음)
\begin{choices}
  \item 매수청구를 한 토지의 소유자는 매수의무자가 그 토지를 매수하지 아니하기로 결정한 경우 개발행위허가를 받아 5층 이하의 제1종근린생활시설을 설치할 수 있다.
  \item 매수의무자는 매수하기로 결정한 토지에 대하여 매수 결정을 알린 날부터 2년 이내에 매수하여야 한다.
  \item 도시·군계획시설채권의 상환기간은 10년 이내로 한다.
  \item 지방자치단체인 매수의무자는 토지소유자가 원하는 경우 채권을 발행하여 매수대금을 지급할 수 있다.
  \item 매수의무자는 매수청구를 받은 날부터 6개월 이내에 매수 여부를 결정하여야 한다.
\end{choices}
\end{question}
\answer{1}
